\section{GEMM design and optimization}
\label{sec:gemm}

The ladder exists so the benchmark table attributes each gain to a cause. Table
\ref{tab:ladder} gives the measured result at 4096 cubed; Figure \ref{fig:ladder}
gives every rung at every swept shape, because a single shape hides the fact that
the curve is not flat. The narrative below is the diagnostic log in brief.

\begin{table}[htbp]
\centering
\begin{tabular}{llrr}
\toprule
rung & precision & GFLOP/s & percent of cuBLAS \\
\midrule
naive & fp32 & 1810 & 8.8 \\
shared tiled & fp32 & 1455 & 7.1 \\
register blocked & fp32 & 16378 & 79.4 \\
cp.async double buffered & fp32 & 17047 & 82.7 \\
WMMA & fp16 & 33722 & 58.6 \\
mma.sync PTX & fp16 & 33627 & 58.4 \\
mma.sync + ldmatrix & fp16 & 33729 & 58.6 \\
mma.sync + ldmatrix + swizzle & fp16 & 54602 & 94.8 \\
CUTLASS reference line & fp16 & 58669 & 101.9 \\
WMMA (BF16) & bf16 & 31148 & 54.1 \\
\bottomrule
\end{tabular}

\caption{GEMM ladder at 4096 cubed, each rung versus cuBLAS of the same precision,
measured on the RTX 5070.}
\label{tab:ladder}
\end{table}

\begin{figure}[htbp]
\centering
\includegraphics[width=0.95\textwidth]{ladder.pdf}
\caption{The ladder as percent of cuBLAS against shape, every rung, every swept
size. Percent is of cuBLAS \emph{at the same precision}: the FP32 rungs are
measured against cuBLAS SGEMM and the tensor rungs against cuBLAS FP16 or BF16,
which is why an FP32 rung can show a higher percent than a faster tensor rung.
\laddercaption}
\label{fig:ladder}
\end{figure}

\subsection{The negative result: shared memory tiling loses to naive}
\label{sec:negative}

The first optimization in every CUDA course is to stage the operand tiles through
shared memory. On this part it is a regression. The naive kernel reaches 1908
GFLOP/s at 4096 cubed; the shared memory tiled kernel that was supposed to improve
on it reaches 1584, which is 17 percent slower. This is the most interesting
result in the project, and it is reported here as a section rather than as an
aside because it inverts what the textbook says to do first.

Round 1 of the diagnostic log names the mechanism, and it has two halves.

The first half is what tiling costs. The tiled kernel uses a 32 by 32 block, which
is 1024 threads, so exactly one block fits per SM: about 66.7 percent occupancy
against the naive kernel's full occupancy. Its top warp stall reason is the MIO
queue, the shared memory pipe. Tiling bought a shared memory bottleneck and paid
for it with residency.

The second half is what tiling was supposed to buy, and did not. The naive kernel
is not bandwidth bound in the first place. Its L1 hit rate is 87.5 percent against
the tiled kernel's 66.7, and its DRAM throughput is only 37 percent, because the
cache hierarchy on this part absorbs the redundant reads that tiling exists to
remove. The redundant traffic never reaches DRAM, so eliminating it wins nothing,
and the shared memory staging is pure added cost.

Two things follow. The technique is not wrong; the ordering is. What makes the
tiled decomposition pay is register blocking on top of it, which is the next rung
and a 5.5 times gain. And the finding is a single GPU finding. The mechanism is a
statement about the ratio of cache capacity to working set, so it should invert on
a part with a smaller last level cache, and the study that would establish that
needs three GPUs of different generations. I do not have them. Section
\ref{sec:limitations} names that as a limitation rather than letting the sentence
generalize on its own.

\subsection{The rest of the ladder}

\paragraph{Register blocking (Round 2).} Giving each thread an 8 by 8 register tile
lifts throughput to 10.4 TFLOP/s, 5.5 times naive and the decisive step. Warp
cycles per issued instruction fall from 45 to 9. The new limiter is the shared read
pipe and a 33 percent occupancy ceiling from 106 registers per thread.

\paragraph{cp.async double buffering (Round 3).} Overlapping the next tile's global
load with the current tile's math takes the FP32 kernel to 16.2 TFLOP/s, 73 percent
of cuBLAS. The FP32 path is now occupancy limited (149 registers per thread, 17
percent occupancy) and is accepted as the top FP32 variant, its gap diagnosed.

\paragraph{Tensor cores (Rounds 4 and 5).} Moving the inner product onto the tensor
cores with WMMA doubles throughput to 36 TFLOP/s, but the cores are starved. In the
Round 4 detail page (\texttt{experiments/results/ncu/round04/wmma\_fp16\_4096.txt})
the GPU Speed Of Light section reads \texttt{Compute (SM) Throughput} 53.72 percent
against \texttt{Memory Throughput} 92.40 percent, and the row that actually names
the culprit is \texttt{L1/TEX Cache Throughput} at 94.38 percent: the fragment loads
saturate the shared read path while the tensor pipe idles. Dropping to mma.sync PTX
with hand placed fragments changes nothing, which
is the informative result: the instruction path was never the bottleneck, the
fragment loads are.

\paragraph{Toward the gate (Rounds 6 to 9).} ldmatrix on the small tile does not
move the number either, because at a 64 by 64 tile the bytes to compute ratio is
too high. A 128 by 128 tile with cp.async double buffering takes it to 80 percent
of cuBLAS and flips the Speed of Light to co limited. A three stage pipeline is
tried and reverted (it costs occupancy). The bank conflict counter then shows the
real problem: 219 million shared load conflicts. A shared memory swizzle, applied
identically to the stores and the loads so correctness is unconditional, cuts them
to 33.7 million, collapses the memory Speed of Light from 76 to 30 percent, and
lifts compute to 83 percent. The kernel now reaches 89.9 percent of cuBLAS at 4096
cubed and 88.9 percent at 8192 cubed, and passes the compute bound gate.

\subsection{The remaining gap, decomposed}
\label{sec:gap}

It is not a flat 90 percent across the range, and the shape of the shortfall is
itself evidence. Table \ref{tab:gap} splits the gap into the part that wave
quantization explains and the part that it does not. The quantization column is
arithmetic, not a measurement: at a 128 by 128 tile the output is
$\lceil m/128 \rceil \times \lceil n/128 \rceil$ tiles competing for 96 resident
block slots (48 SMs, two blocks each, which is what the shared memory budget
allows), and the loss is the fraction of the last wave that runs empty.

\begin{table}[htbp]
\centering
\begin{tabular}{lrrlrr}
\toprule
shape & GFLOP/s & percent of cuBLAS & CTAs into 96 slots & quantization loss & residual \\
\midrule
1024 cubed & 32,303 & 63.5 & 64, 0.67 waves & 33.3 & 3.2 \\
2048 cubed & 51,493 & 80.5 & 256, 2.67 waves & 11.1 & 8.4 \\
4096 cubed & 54,985 & 89.9 & 1024, 10.67 waves & 3.0 & 7.0 \\
8192 cubed & 58,060 & 88.9 & 4096, 42.67 waves & 0.8 & 10.3 \\
\bottomrule
\end{tabular}

\caption{The gap to cuBLAS decomposed. Throughput and percent are measured; the
wave and quantization columns are computed from the tile geometry; the residual is
what quantization does not account for.}
\label{tab:gap}
\end{table}

Quantization dominates at 1024 cubed and is a minority term above 2048. The
residual at 4096 and 8192 is a different gap, and my working hypothesis is that
most of it is kernel selection: cuBLAS carries a family of tile shapes and chooses
per shape, while the single rung above optimizes one fixed 128 by 128 by 32 tile.
That hypothesis is what Section \ref{sec:mechanisms} exists to test, and it is
labelled a hypothesis until the sweep tests it.
